% Chapter 1: Introduction: The Paradigm Shift to Infinite-Dimensional Learning
% Part I of the book.

\chapter{Introduction: The Paradigm Shift to Infinite-Dimensional Learning}
\label{ch:part1-introduction-the-paradigm-shift-to-infinit}
\label{sec:intro}
% -----------------------------------------------------------------------------

Numerical simulation answers one question at a time: given a permeability field, an
initial state, or boundary data, a
\index{finite volume method (FVM)}\emph{finite volume method} (FVM) or a
\index{finite element method (FEM)}\emph{finite element method} (FEM) returns the matching
solution field, and changing one coefficient carries nothing over from the previous solve.
The workloads that motivate this book do not pose single questions. Shape optimization
queries the same governing equations for thousands of candidate geometries, ensemble weather
forecasting integrates a perturbed initial condition many times, and reservoir risk
assessment samples a distribution over subsurface heterogeneity; every setting in Part~III
asks for many solutions of one operator. We call this the
\index{multi-query surrogate problem}\index{operator learning}\emph{multi-query surrogate problem}: the quantity of interest is not a solution but the map from data to solution, and
it must be evaluated at whatever resolution each query requires.

Stretching a conventional network over such a workload does not remove the mismatch, because
these architectures are maps between fixed-dimensional Euclidean spaces: parameter count,
weight-tensor shapes, and the physical meaning of each activation are all set by the number
of grid points present at training time. The alternative developed in this book takes the
\emph{operator} as the object of learning---a map defined on function spaces, written down
without reference to a mesh, then evaluated on whichever discretization a query supplies.
Parts~II and~III inherit that change of target, and it is not a change of vocabulary: it
raises three questions finite-dimensional learning theory never asks. What spaces does the
network consume and produce, and how much regularity do their elements carry? What must the
target map satisfy for approximation to be possible at all? And what does it mean for a
machine to apply a map not defined on a finite list of numbers?

This chapter states the problem and the shape of its answers; it proves almost
nothing. The arguments come from the rest of Part~I.
\chapref{ch:part1-function-spaces-duality-and-sobolev-regula} fixes the spaces, the norms,
and the compactness machinery;
\chapref{ch:part1-partial-differential-equations-as-solution} verifies that the solution maps
of Darcy flow, Burgers' equation, and Navier--Stokes carry the continuity and smoothing
properties the theory demands;
\chapref{ch:part1-operator-approximation-theory-and-universa} supplies the density results;
and \chapref{ch:part1-the-fourier-neural-operator-spectral-theor} with
\chapref{ch:part1-discretization-invariance-and-continuum-li} deliver an architecture and a
mesh-independence statement. The reader should bring measure-based probability, a first
course in functional analysis, Sobolev spaces at the level of
\cite{brezis2011functional, evans2010partial}, the discrete Fourier transform, and enough
PyTorch to read a layer definition; results Part~I invokes without proof are collected in
Appendix~\ref{app:a}.

\section{The Multi-Query Bottleneck in Simulation-Based Science}

Classical deep learning architectures---multilayer perceptrons
\index{multilayer perceptron (MLP)}(MLPs), convolutional neural networks
\index{convolutional neural network (CNN)}(CNNs), recurrent neural networks
\index{recurrent neural network (RNN)}(RNNs)---are formulated as maps between
finite-dimensional Euclidean spaces,
\begin{equation}
  \label{eq:ch01-euclidean-map}
  f_\theta : \mathbb{R}^{d_{\text{in}}} \longrightarrow \mathbb{R}^{d_{\text{out}}}, \qquad \theta \in \Theta \subset \mathbb{R}^p,
\end{equation}
with $p$ trainable parameters and with the input and output dimensions $d_{\text{in}},
d_{\text{out}}$ fixed by the architecture before any data are seen. Equation
\eqref{eq:ch01-euclidean-map} is not one design's limitation but the shared contract of the
class: a tensor of size $d_{\text{in}}$ must arrive with its entries in a prescribed order.

Applied to systems governed by \index{partial differential equation (PDE)}\emph{partial
differential equations} (PDEs), such a model therefore consumes discrete approximations of
continuous fields, sampled at fixed spatial or temporal locations $\{x_i\}_{i=1}^{N}
\subset \Omega$: an input field $a$ enters as the vector $\big(a(x_1), \dots, a(x_N)\big)
\in \mathbb{R}^{N d_a}$, so that $d_{\text{in}} = N d_a$, with $d_a$ the number of components
of $a$ and $d_u$ the number of components of the output. The consequence is that the
learned parameters $\theta$ remain intrinsically coupled to the \index{mesh resolution!coupling of network parameters to}mesh resolution $N$, the grid
topology, and the sampling geometry used during training. Refine the resolution
($N \to N'$), deform the mesh, or move the sensors, and the finite-dimensional model becomes
inapplicable without ad hoc interpolation, retraining, or loss of fidelity. The usual repair,
re-interpolation, is the least defensible: resampling a surrogate's output onto $N'$ points
introduces an error training never saw or bounded, so the reported accuracy describes grid
$N$, not the field.

The economics define what the machinery of this book is for. Let $M$ be the number of queries, let $c_{\mathrm{sim}}(N)$ be the cost of one classical
solve at $N$ unknowns, let $c_{\mathrm{train}}$ be the one-time cost of building a surrogate,
and let $c_{\mathrm{eval}}(N)$ be the cost of one surrogate evaluation:
\begin{equation}
  \label{eq:ch01-cost-bookkeeping}
  C_{\mathrm{sim}}(M, N) = M\, c_{\mathrm{sim}}(N),
  \qquad
  C_{\mathrm{sur}}(M, N) = c_{\mathrm{train}} + M\, c_{\mathrm{eval}}(N).
\end{equation}
Equation \eqref{eq:ch01-cost-bookkeeping} is arithmetic, not a theorem, and it invites three
qualifications. First, the surrogate wins only when $c_{\mathrm{eval}}(N) <
c_{\mathrm{sim}}(N)$, with crossover at $M^{\star} = c_{\mathrm{train}} /
\big(c_{\mathrm{sim}}(N) - c_{\mathrm{eval}}(N)\big)$; a single design check should still be
handed to a classical solver \cite{leveque2002finite, quarteroni2008numerical}. Second,
$c_{\mathrm{train}}$ is not free of the solver it replaces, because supervised training
produces its paired samples by running that solver repeatedly; amortization presupposes data
generation as a sunk cost shared across many downstream queries.
Third, and most important for what follows, $c_{\mathrm{eval}}(N)$ still grows with $N$. A
spectral forward pass costs $\mathcal{O}(N \log N)$ operations on a uniform periodic grid,
whereas a branch-and-trunk network costs an amount set by its widths and by the number of
sensor and query points rather than by the evaluation grid
\cite{li2021fourier, lu2021learning}. The gain of operator learning is amortization of a
family of solves, not removal of the resolution dependence.

What the fixed-mesh formulation actually forfeits is not speed but scope. Three distinct
failures are entangled in \eqref{eq:ch01-euclidean-map} once the mesh changes: resolution
coupling, since weights sized for $N$ nodes have no defined action on $N'$ nodes; geometry
coupling, since convolutional weights are indexed by lattice offsets, so a refinement shrinks
their physical footprint---the argument made quantitatively in
\chapref{ch:part1-discretization-invariance-and-continuum-li}; and a semantic failure often missed in
benchmarks, namely that a network trained on one sampling of one field learns \emph{a}
solution rather than the rule that generates solutions from data.
The first two can be patched by re-interpolation at some cost in accuracy; the third cannot,
because the network represents an object the application does not need. Diagnosing
it requires a target defined independently of any discretization, which is the next
section's subject.

\section{Solution Operators as Targets of Learning}

Scientific and physical laws operate on continuous fields. Velocity fields satisfy
$\mathbf{u}(\cdot, t) \in (L^{2}(\Omega))^{d}$, temperature distributions satisfy
$T(\cdot) \in H^{1}(\Omega)$, and permeability tensors satisfy
$\mathbf{K}(\cdot) \in L^{\infty}(\Omega; \mathbb{R}_{\text{sym}}^{d \times d})$, where
$\Omega \subseteq \mathbb{R}^{d}$ is the spatial domain, $L^{p}(\Omega)$ denotes functions
whose $p$-th power is integrable, and $H^{s}(\Omega)$ denotes the \index{Sobolev space}Sobolev space of elements
with $s$ derivatives in $L^{2}(\Omega)$; \chapref{ch:part1-function-spaces-duality-and-sobolev-regula}
develops both. Each of these is an element of an infinite-dimensional vector space, and the
map that carries an initial state, a boundary forcing, or a coefficient field to the system
state at time $t$ is a \textbf{solution operator}
\index{solution operator}
\begin{equation}
  \label{eq:ch01-solution-operator}
  \mathcal{G} : \mathcal{A} \longrightarrow \mathcal{U},
\end{equation}
where the \index{A@$\mathcal{A}$ (input function space)}input space $\mathcal{A}$ and the
\index{U@$\mathcal{U}$ (output function space)}output space $\mathcal{U}$ are infinite-dimensional Banach or Hilbert function
spaces over $\Omega$
\cite{kovachki2023neural, stuart2010inverse}. Equation \eqref{eq:ch01-solution-operator} is
deliberately terse, and the two spaces are usually different: for the Darcy problem treated
in \chapref{ch:part1-partial-differential-equations-as-solution}, the coefficient enters from
$L^{\infty}(\Omega)$ while the solution lands in $H^{1}_{0}(\Omega)$. The same notation covers
the time-$t$ map $u_{0} \mapsto u(\cdot, t)$, the one-step map
$u(t) \mapsto u(t + \Delta t)$ used for autoregressive rollouts, and parameter-to-state maps,
in which finitely many physical parameters enter as functions on $\Omega$---which is why the
target is stated as \eqref{eq:ch01-solution-operator} rather than as a
particular PDE.

Contrast this with the \index{Physics-Informed Neural Network (PINN)}physics-informed neural
network (PINN), in which a network approximates
\emph{one} solution by minimizing a residual at collocation points
\cite{raissi2019physics}. A PINN's unknown is a field; the operator learner's unknown is a
map, so a new input case is a forward pass rather than a new optimization. Because the
application's queries are drawn from a distribution, training minimizes
an expected loss under an \index{input measure}input measure
$\mu \in \mathcal{P}(\mathcal{A})$, where $\mathcal{P}(\mathcal{A})$ is the set of Borel
probability measures on $\mathcal{A}$,
\begin{equation}
  \label{eq:ch01-risk}
  \mathcal{R}(\theta)
  = \mathbb{E}_{a \sim \mu}
    \left[ \frac{\big\|\mathcal{G}_{\theta}(a) - \mathcal{G}(a)\big\|_{\mathcal{U}}^{2}}
                {\big\|\mathcal{G}(a)\big\|_{\mathcal{U}}^{2}} \right],
    \qquad a \in \mathcal{A},
\end{equation}
and the normalization in \eqref{eq:ch01-risk} makes errors comparable across resolutions and
across magnitudes---the relative-$L^{2}$ metric used throughout Parts~II and~III. Approximation
theory, however, states its guarantees in a different language, as a worst-case bound $\sup_{a \in K} \|\mathcal{G}_{\theta}(a) -
\mathcal{G}(a)\|_{\mathcal{U}} < \varepsilon$ over a compact set $K \subset \mathcal{A}$. The
two are related but not interchangeable: the supremum controls the risk only when $\mu$ gives
$K$ full measure and $\|\mathcal{G}(a)\|_{\mathcal{U}}$ stays bounded away from zero, whereas
the risk controls the supremum nowhere. The gap between them is the statistical
error that \chapref{ch:part1-operator-approximation-theory-and-universa} and
\chapref{ch:part1-open-theoretical-problems-and-mathematical} return to, and a surrogate that
satisfies one and not the other is a familiar engineering failure.

Turning \eqref{eq:ch01-solution-operator} into something trainable means proposing a
parametric family $\mathcal{G}_{\theta}$, $\theta \in \Theta \subset \mathbb{R}^{p}$, with
three requirements: the parameter space $\Theta$ is defined independently of any
discretization; each $\theta$ determines an operator on the continuum spaces
$\mathcal{A}$ and $\mathcal{U}$; and the resulting map is continuous in the norm of
$\mathcal{U}$. Continuity is not a formality. Networks assemble bounded linear maps and
Lipschitz activations, so every architecture considered here produces continuous
$\mathcal{G}_{\theta}$, while a discontinuous target cannot be approximated uniformly on any
set containing the jump---which is why the density theorems of
\cite{chen1995universal, kovachki2021universal} all place continuity of $\mathcal{G}$ among
their hypotheses. The requirement pushes the burden of proof upstream: before asking whether
the architecture class is rich enough, we must ask whether the solution map is continuous and
on which subsets of $\mathcal{A}$ it is---the content of
\chapref{ch:part1-partial-differential-equations-as-solution}, resting on the regularity
theory of the chapter before it.

\takeaway{The Fundamental Objective of Operator Learning}{Rather than approximating a specific solution vector on a fixed discretization, neural operator theory seeks to parameterize a family of mappings $\mathcal{G}_\theta: \mathcal{A} \to \mathcal{U}$ such that $\mathcal{G}_\theta$ is continuous, mesh-independent, and capable of being queried on arbitrary discretizations without retraining.}

Several realizations of that program recur throughout the book, each trading one kind of
coupling for another. \index{DeepONet}DeepONet (Deep Operator Network) reads the input at $m$ fixed
sensors through a branch network and synthesizes the output at arbitrary query coordinates
through a trunk network, so the output side is mesh-free while the input side stays
sensor-bound \cite{lu2021learning}. The
\index{neural operator!generic, integral-kernel architecture}generic neural operator replaces
sensor vectors with learnable integral kernels acting on whole functions
\cite{kovachki2023neural}, and the \index{Fourier neural operator (FNO)}\emph{Fourier neural
operator} (FNO) restricts those kernels to translation-invariant convolution on periodic
domains, which is what makes the continuum operator cheap to evaluate on a grid
\cite{li2021fourier}. The \index{Graph Neural Operator (GNO)}graph neural operator (GNO) and the
\index{Geometry-Informed Neural Operator (GINO)}geometry-informed neural operator (GINO) move the same machinery onto
unstructured point clouds and surface meshes; the
\index{Physics-Informed Neural Operator (PINO)}physics-informed neural operator (PINO) trains
against weak-form residuals rather than paired data alone \cite{li2023geometry, li2021physics}.
Figure~\ref{fig:architectures} in
\chapref{ch:part1-operator-approximation-theory-and-universa} makes the comparison concrete:
panel~(a) is a CNN or MLP whose parameters are bound to one discretization, panel~(b) a
branch-and-trunk network that queries continuous output coordinates through fixed input
sensors, panel~(c) a spectral network that keeps operator semantics under mesh change; and
Table~\ref{tab:comparison} of
\chapref{ch:part1-comparative-theoretical-taxonomy-of-archit} records the same differences
property by property. Part~II holds the implementations, in PyTorch \cite{paszke2019pytorch},
and further descendants---the U-FNO (U-Net-enhanced Fourier neural operator), which interleaves
multi-scale residual blocks with spectral convolution for multiphase flow, and the
\index{adaptive Fourier neural operator (AFNO)}adaptive Fourier neural operator (AFNO) for
global weather---are analyzed
in \chapref{ch:part1-comparative-theoretical-taxonomy-of-archit}, indexed in
Appendix~\ref{app:e}, and applied in Part~III
\cite{wen2022ufno, pathak2022fourcastnet}.

\section{What Changes When the Dimension Is Not Fixed}

Once $\mathcal{A}$ and $\mathcal{U}$ are infinite-dimensional, three intuitions that carry a
finite-dimensional learning theory stop working, and the rest of Part~I exists largely to
replace them. The first is compactness. In $\mathbb{R}^{d}$ a closed and bounded set is
compact, so ``approximate uniformly on the data range'' is an ordinary instruction. In
infinite dimensions the closed unit ball is never compact: bounded sequences need not admit
convergent subsequences at all. Uniform approximation statements must therefore be restricted
to compact subsets $K \subset \mathcal{A}$, and compactness in function space has to be
purchased with regularity rather than assumed. The mechanism is the \index{Rellich--Kondrachov compact embedding}Rellich--Kondrachov
compact embedding, which makes bounded sets in $H^{s}(\Omega)$ precompact in
$H^{t}(\Omega)$ for $s > t \ge 0$; it is proved and put to work in
\chapref{ch:part1-function-spaces-duality-and-sobolev-regula} and is the reason the \index{universal approximation!of operators}density
theorems in \chapref{ch:part1-operator-approximation-theory-and-universa} are stated on
compacta. \index{compactness!failure in infinite dimensions}

The second replacement concerns smoothing, which turns out to be the property that makes an
operator learnable at all. Viscous parabolic problems improve their inputs: for viscous
Burgers' equation with $\nu > 0$, initial data $u_{0} \in L^{2}(\mathbb{T})$ evolve to
$u(\cdot, t) \in H^{s}(\mathbb{T})$ for every $s > 0$. Elliptic problems gain regularity by a
fixed amount: for smooth Darcy coefficients, classical estimates turn $L^{2}$ forcing into an
$H^{2}$ solution, as
\chapref{ch:part1-partial-differential-equations-as-solution} makes explicit. Where the
target behaves as
$\mathcal{G} : H^{s} \to H^{s+\alpha}$ with $\alpha > 0$, the image of a bounded set is
precompact and truncating its range to finitely many Fourier modes costs an amount that
decays like $k_{\max}^{-\alpha}$, which is the content of the spectral convergence result in
\chapref{ch:part1-the-fourier-neural-operator-spectral-theor} \cite{kovachki2021universal}.
Where smoothing fails---inviscid hyperbolic limits, shock formation, transport of
discontinuous data---regularity is destroyed rather than created, and neither the compactness
nor the rate argument survives without new hypotheses. That asymmetry, not a shortage of
layers, is where the honest limits of the theory sit;
\chapref{ch:part1-open-theoretical-problems-and-mathematical} lists it as open.

The third replacement is that a single ``error'' becomes a sum of errors that behave
differently under refinement. Let $\theta^{\star}$ denote a parameter minimizing
$\mathcal{R}$ in \eqref{eq:ch01-risk} within the chosen class, let $\widehat{\theta}$ denote
the parameter produced from finitely many paired samples, and let $\mathcal{G}_{\theta, N}$
denote the discretely implemented operator, with integrals replaced by quadrature on $N$
nodes. Adding and subtracting terms gives
\begin{equation}
  \label{eq:ch01-error-split}
  \begin{aligned}
    \big\|\mathcal{G}(a) - \mathcal{G}_{\widehat{\theta}, N}(a)\big\|_{\mathcal{U}}
    \;\leq\;&
    \underbrace{\big\|\mathcal{G}(a) - \mathcal{G}_{\theta^{\star}}(a)\big\|_{\mathcal{U}}}_{\text{approximation}}
    \\
    &+ \underbrace{\big\|\mathcal{G}_{\theta^{\star}}(a) - \mathcal{G}_{\theta^{\star}, N}(a)\big\|_{\mathcal{U}}}_{\text{discretization}}
    + \underbrace{\big\|\mathcal{G}_{\theta^{\star}, N}(a) - \mathcal{G}_{\widehat{\theta}, N}(a)\big\|_{\mathcal{U}}}_{\text{statistical}} .
  \end{aligned}
\end{equation}
The decomposition \eqref{eq:ch01-error-split} follows from the triangle inequality alone, so
its value lies in which terms the architecture makes separately controllable; authors draw
the cut differently. For DeepONet the natural split is sensor count, trunk basis dimension,
and network error \cite{lanthaler2022error}; for spectral operators, mode truncation appears
as its own term alongside discretization and statistical error
\cite{subedi2024controlling}, the three-term form written as \eqref{eq:ch06-decomp} in
\chapref{ch:part1-discretization-invariance-and-continuum-li}. A fourth term is missing in
practice: $\widehat{\theta}$ comes from stochastic gradient descent rather than from an exact
minimizer, so optimization error sits underneath \eqref{eq:ch01-error-split} and is treated
empirically in Part~II. What the split does establish is a warning about reporting: a
``relative error'' measured at the training resolution confounds the approximation term with
the discretization term, and if the second is not small the reported number describes the
quadrature rather than the learned operator.
\index{error decomposition!approximation, discretization, and statistical}
\index{smoothing property!of parabolic and elliptic problems}

Notation follows the front matter rather than this chapter: Table~\ref{tab:notation-master}
fixes $\Omega$, $L^{p}(\Omega;\mathbb{R}^{d_v})$,
$H^{s}(\Omega)=W^{s,2}(\Omega)$, $C^{k}(\overline{\Omega})$, the generic spaces $X, Y$,
the operator space $\mathcal{L}(X,Y)$, the duality pairing
$\langle \cdot, \cdot \rangle_{X^{*},X}$, the transform pair
$\mathcal{F}, \mathcal{F}^{-1}$, the \index{Td@$\mathbb{T}^{d}$ (periodic torus)}torus $\mathbb{T}^{d} \cong [0,1)^{d}$, and weak
derivatives $D^{\alpha}u$; the glosses above are reminders, and the table governs where
they differ. Beyond it,
this book reserves $N$ for the number of grid points, $M$ for the number of queries, and
$d_a$ and $d_u$ for input and output channel dimensions, while $p$ denotes the number of
trainable parameters only outside a space name---in $L^{p}$ and $W^{k,p}$ it is always the
integrability exponent. Thus $\mathcal{A}$ and $\mathcal{U}$ name spaces while
$\mathcal{G}$ and $\mathcal{G}_{\theta}$ name maps, and Part~II re-encodes the same
fields as tensors of shape $(B, C, H, W)$, where $B$ indexes the batch and $C$ the channels
$d_a$ or $d_u$; \chapref{ch:part2-the-practitioner-s-paradigm-infinite-dimen} makes that
translation explicit.
\index{notation!master table}

Finally, a word on what operator learning does not buy. A surrogate inherits the support of the
input measure $\mu$: an error measured on smooth, well-resolved coefficient fields licenses no claim
about data from outside that support, and the framework's guarantees are stated relative to a
chosen space and a chosen compact subset of it
\cite{kovachki2023neural}. Sample-complexity bounds under a general measure on function space
are not yet sharp, a gap argued in
\chapref{ch:part1-open-theoretical-problems-and-mathematical}. Conserved quantities and
boundary conditions are not enforced by \eqref{eq:ch01-risk} unless the architecture or the
loss imposes them, a matter taken up in Part~III. And a spectral representation of
shock-dominated data exhibits Gibbs oscillations near the jump
\index{Gibbs phenomenon}, so smoothness-driven accuracy and discontinuity resolution are in
tension---one of the open problems closing Part~I. Stating these limits first means the
theorems that follow read as conditional results, hypotheses attached.

\bigskip
\noindent What this chapter has asserted and not proved is the schedule for the rest of
Part~I. We demanded that $\mathcal{A}$ and $\mathcal{U}$ be complete spaces in which norms
measure more than point values, that the target map be continuous, that the input set be
compact, and that quadrature on $N$ nodes converge to the continuum integral as $N$ grows; we
justified none of it. \chapref{ch:part1-function-spaces-duality-and-sobolev-regula} supplies
the infrastructure: Banach and Hilbert structure and the Riesz representation theorem, which
identifies the dual of a Hilbert space with the space itself and thereby makes sense of
pairings such as $\langle f, v \rangle_{H^{-1}, H^{1}_{0}}$ used throughout Part~I; weak
derivatives and the Sobolev spaces $W^{k,p}$ and $H^{s}$, including the Fourier
characterization on $\mathbb{T}^{d}$ that the spectral analysis of
\chapref{ch:part1-the-fourier-neural-operator-spectral-theor} will consume; the
Rellich--Kondrachov and Sobolev embedding theorems, which convert regularity into the
compactness on which the approximation term of \eqref{eq:ch01-error-split} rests; and Bochner spaces
$L^{p}(0,T;X)$, without which an evolution operator has no topology of its own. Readers who command
these tools may treat the next chapter as a review and move to the PDE examples; those who do
not should not skip it, because every constant and hypothesis in Part~I is stated in its
language.
